\section{Kiến thức chuẩn bị}

\subsection{Giải tích}

Ngoại trừ định nghĩa về không gian Hilbert, các định nghĩa sau được tham khảo
từ \cite[chương~2, 3 và~7]{Rudin1976principles}.

\begin{defn}[Không gian \tr{metric}]
  Một tập $X$, với các phần tử gọi là điểm, được gọi là không gian \tr{metric}
  (metric space) nếu với hai điểm $p$ và $q$ bất kì của $X$, tồn tại một số
  thực $d(p, q)$ tương ứng, được gọi là khoảng cách từ $p$ đến $q$, sao cho
  \begin{itemize}
    \item $d(p, q) > 0$ nếu $p \ne q$; $d(p, p) = 0$;
    \item $d(p, q) = d(q, p)$;
    \item $d(p, q) \le d(p, r) + d(r, q)$ với mọi $r \in X$.
  \end{itemize} Mọi hàm thỏa cả ba tính chất trên được gọi là hàm khoảng cách
  (distance function) hay là \tr{metric} (metric).
\end{defn}

\begin{defn}
  Cho không gian \tr{metric} $X$, tất cả các điểm và tập sau đây đều xem là
  điểm hoặc tập con của $X$. \begin{itemize}
    \item Một lân cận (neighborhood) của $p$ là một tập $N_r(p)$ gồm tất cả các
      điểm $q$ sao cho $d(p, q) < r$ với $r > 0$ nào đó. $N_r(p)$ còn được gọi
      là quả cầu (ball) mở có tâm là $p$ và bán kính $r$.
    \item Một điểm $p$ được gọi là điểm giới hạn (limit point) của tập $E$ nếu
      mọi lân cận của $p$ đều chứa ít nhất một điểm $q \ne p$ thuộc $E$.
    \item Điểm $p$ được gọi là điểm cô lập nếu nó không phải điểm giới hạn.
    \item Tập $E$ được gọi là đóng (closed) nếu mọi điểm giới hạn của $E$ đều
      thuộc $E$.
    \item Một điểm $p$ được gọi điểm trong (interior point) của $E$ nếu tồn tại
      một lân cận của $p$ là tập con của $E$.
    \item Tập $E$ được gọi là mở (open) nếu mọi điểm của $E$ đều là điểm trong
      của $E$.
    \item Tập $E$ được gọi là hoàn hảo (perfect) nếu $E$ đóng và mọi điểm của
      $E$ đều là điểm giới hạn của $E$.
    \item Tập $E$ được gọi là bị chặn (bounded) nếu tồn tại số thực $M$ và điểm
      $q$ sao cho $d(p, q) < M$ với mọi $p \in E$.
    \item Tập $E$ được gọi là trù mật (dense) trong $X$ nếu mọi điểm của $X$
      đều là điểm giới hạn của $E$ hay thuộc $E$.
  \end{itemize}
\end{defn}

\begin{defn}[Phủ mở]
  Một phủ mở (open cover) của một tập $E$ trong không gian \tr{metric} $X$ là
  một họ $\{G_\alpha\}$ các tập con mở của $X$ sao cho $E \subseteq \cup_\alpha
  G_\alpha$.
\end{defn}

\begin{defn}[Tập \tr{compact}]
  Tập con $K$ của không gian \tr{metric} $X$ được gọi là \tr{compact} (compact)
  nếu mọi phủ mở của $K$ đều tồn tại một phủ mở con hữu hạn trong đó. Rõ ràng
  hơn, nếu ${G_\alpha}$ là một phủ mở của $K$ thì tồn tại các chỉ số hữu hạn
  $\alpha_1$, $\alpha_2$, $\ldots\,$, $\alpha_n$ sao cho $\{G_{\alpha_i} \colon
  i \in [n]\}$ là một phủ mở của $K$.
\end{defn}

\begin{defn}[Không gian \tr{compact} địa phương]
  Không gian \tr{metric} $X$ được gọi là \tr{compact} địa phương (locally
  compact) nếu mọi điểm $p \in X$ đều có một lân cận \tr{compact}, tức là tồn
  tại một tập mở $U$ và một tập \tr{compact} $K$ sao cho $ x\in U \subseteq K$.
\end{defn}

\begin{defn}[Dãy hội tụ]
  Một dãy $(p_n)$ trong không gian \tr{metric} $X$ được gọi là hội tụ
  (converge) nếu tồn tại một điểm $p \in X$ thỏa tính chất sau: Với mọi
  $\varepsilon > 0$, tồn tại số nguyên $N$ sao cho $n \ge N$ kéo theo $d(p,
  p_n) \le \varepsilon$. (Ở đây $d$ kí hiệu khoảng cách trong $X$.) Và ta kí
  hiệu \begin{equation}
    \lim_{n\to\infty} p_n = p.
  \end{equation}
\end{defn}

\begin{defn}[Dãy Cauchy]
  Một dãy $(p_n)$ trong không gian \tr{metric} $X$ được gọi là dãy Cauchy
  (Cauchy sequence) nếu với mọi $\varepsilon > 0$, tồn tại số nguyên $N$ sao
  cho $d(p_n, p_m) < \varepsilon$ nếu $n \ge N$ và $m \ge N$.
\end{defn}

\begin{defn}[Không gian \tr{metric} đầy đủ]
  Không gian \tr{metric} được gọi là đầy đủ (complete) nếu mọi dãy Cauchy trong
  đó đều hội tụ.
\end{defn}

\begin{defn}[Hội tụ điểm]
  Giả sử $(f_n)_{n = 1}^\infty$ là dãy các hàm được định nghĩa trên tập $E$, và
  giả sử mọi dãy $\bigl( f_n(x) \bigr )$ đều hội tụ với mọi $x \in E$. Khi đó,
  ta có thể định nghĩa hàm $f$ bởi \begin{equation}
    f(x) = \lim_{n\to \infty} f_n(x) \quad (x \in E).
  \end{equation} Khi này $f$ được gọi hội tụ điểm  (pointwise limit) của dãy
  hàm $(f_n)_{n = 1}^\infty$ trên $E$.
\end{defn}

\begin{defn}[Hội tụ đều]
  Ta nói dãy các hàm $(f_n)_{n = 1}^\infty$ hội tụ đều (converge uniformly)
  trên $E$ đến một hàm $f$ nếu với mọi $\varepsilon > 0$, tồn tại số nguyên $N$
  sao cho $n \ge N$ kéo theo $\bigl \lvert f_n(x) - f(x) \bigr \rvert \le
  \varepsilon$ với mọi $x \in E$.
\end{defn}

\begin{defn}[Tiên đề về tích trong]
  Giả sử $V$ là một không gian \tr{vector} trên trường số vô hướng $\mathbb{F}$
  (trường thực hoặc phức). Ánh xạ $\langle \cdot, \cdot \rangle \colon V \times
  V \to \mathbb{F}$ được gọi là tích trong trên $V$ nếu nó thỏa mãn ba tiên đề
  sau với mọi $x, y , z \in V$ và mọi $\alpha, \beta \mathbb{F}$:
  \begin{itemize}
    \item Đối xứng liên hợp (conjugate symmetry): $\langle x, y \rangle
      = \overline{\langle y, x \rangle}$.
    \item Tuyến tính theo biến thứ nhất (linearity in the first argument):
      $\langle \alpha x + \beta y, z \rangle = \alpha\langle x, z \rangle +
      \beta \langle y, z \rangle$.
    \item Xác định dương (positive-definiteness): $\langle x, x \rangle \ge 0$
      và $\langle x, x \rangle  = 0$ khi và chỉ khi $x = 0$.
  \end{itemize}
\end{defn}

\begin{defn}[Không gian Hilbert]
  Không gian Hilbert (Hilbert space) là một không gian \tr{vector} được gắn với
  một tích trong $\langle \cdot, \cdot \rangle$ và đầy đủ với khoảng cách được
  sinh ra từ chuẩn được định nghĩa như sau: \begin{equation}
    \forall x \in \mathbb{H}, \quad \lVert x \rVert_\mathbb{H} = \sqrt{\langle
    x, x \rangle }.
  \end{equation}
\end{defn}

\subsection{Ngôn ngữ hình thức}

Các định nghĩa sau được tham khảo từ \cite[chương~1]{Phuong2018}.

\begin{defn}[Bảng chữ cái]
  Bảng chữ cái là tập hữu hạn khác rỗng các kí tự (hay kí hiệu) thường được
  biểu diễn bởi chữ cái Hy Lạp $\Sigma$.
\end{defn}

\begin{defn}[Chuỗi]
  Một chuỗi là dãy hữu hạn các kí tự thuộc bảng chữ cái $\Sigma$. Đặc biệt,
  chuỗi rỗng, kí hiệu là $\epsilon$, có thể hình từ mọi bảng chữ cái $\Sigma$
  bất kì.
\end{defn}

\begin{defn}[Ngôn ngữ]
  Cho bảng chữ cái $\Sigma$, ngôn ngữ lạ tập hợp $\mathcal{L}$ của các chuỗi
  được hình thành từ $\Sigma$. Trong ngữ cảnh này, một chuỗi của $\mathcal{L}$
  được gọi là câu.
\end{defn}

\begin{defn}[Văn phạm]
  Văn phạm $\mathcal{G}$ là bộ bốn thành phần \begin{equation}
    \mathcal{G} = (V, T, P, S),
  \end{equation} trong đó \begin{itemize}
    \item $V$ là tập hữu hạn các kí hiệu không kết thúc;
    \item $T$ là tâp hữu hạn các kí hiệu kết thúc;
    \item $P$ là tập hữu hạn các luật sinh của văn phạm;
    \item $S \in V$ là kí hiệu bắt đầu của văn phạm.
  \end{itemize}
\end{defn}

\begin{defn}[Ngôn ngữ của văn phạm]
  Cho văn phạm $\mathcal{G}$, ngôn ngữ $\mathcal{L}(\mathcal{G})$ của văn phạm
  là tập tất cả các câu mà một văn phạm có thể sinh ra được.
\end{defn}

\begin{defn}[\tr{Automaton}]
  Một \tr{automaton} gồm (không nhất thiết đầy đủ ngoại trừ bộ điều khiển) các
  thành phần:\begin{itemize}
    \item Tệp nhập: Dãy tuyến tính vô hạn các ô nhớ chứa thông tin cần xử lí.
      Mỗi đơn vị thông tin được hình thức hóa dưới dạng một kí hiệu thuộc bảng
      chữ cái, lưu trữ trên một ô nhớ.
    \item Tệp xuất: Dãy tuyến tính vô hạn các ô nhớ chứa thông tin sau xử lí.
    \item Tệp nhớ: Dãy tuyến tính vô hạn các ô nhớ lưu trữ thông tin trong quá
      trình vận hành. Một đơn vị thông tin được hình thức hóa dưới dạng một kí
      hiệu (không nhất thiết thuộc bảng chữ cái), lưu trữ trên một ô nhớ.
    \item Bộ điều khiển: Tập hữu hạn các trạng thái điều khiển.
  \end{itemize} Các trạng thái điều khiển bao gồm: \begin{itemize}
    \item Trạng thái kết thúc (hay chấp nhận): có ít nhất một trạng thái kiểu
      này;
    \item Trạng thái không kết thúc: các trạng thái còn lại;
    \item Trạng thái bắt đầu: là trạng thái đặc biệt đại diện cho nơi máy ảo
      ``đứng'' trước khi vận hạnh, có thể thuộc một trong hai loại trên.
  \end{itemize} Các trạng thái được kết nối với nhau qua hàm truyền. Tham số
  đầu vào của hàm là những thông tin khả dụng đối với \tr{automaton} tại thời
  điểm đang xét, thường là trạng thái hiện hành, kí tự đang được đọc trên tệp
  nhập và có thể cả kí hiệu đang được đọc trên tệp nhớ. Đầu ra của hàm truyền
  chính là trạng thái kế tiếp mà \tr{automaton} sẽ chuyển đến, đồng thời, nội
  dung của ô nhớ hiện tại trên các tệp có thể được duy trì hoặc thay đổi.
\end{defn}

\begin{defn}[Ngôn ngữ của \tr{automaton}]
  Ngôn ngữ của \tr{automaton} là tập tất cả các chuỗi nhập mà nó chấp nhận, tức
  là sau khi xử lý, kết quả ở trạng thái chấp nhận.
\end{defn}
